\documentclass[11pt]{article}

\usepackage[margin=1in]{geometry}
\usepackage{setspace}
\usepackage{amsmath,amssymb}
\usepackage{booktabs}
\usepackage{multirow}
\usepackage{graphicx}
\usepackage{subcaption}
\usepackage{float}
\usepackage{longtable}
\usepackage{array}
\usepackage{xcolor}
\usepackage[plainpages=false,pdfpagelabels]{hyperref}
\usepackage{url}
\usepackage{enumitem}
\usepackage{microtype}

\hypersetup{
  colorlinks=true,
  linkcolor=blue,
  citecolor=blue,
  urlcolor=blue
}

\setlength{\parskip}{0.25em}
\setlength{\parindent}{1.2em}
\setstretch{1.15}
\setlength{\textfloatsep}{10pt plus 2pt minus 2pt}
\setlength{\floatsep}{8pt plus 2pt minus 2pt}
\setlength{\intextsep}{8pt plus 2pt minus 2pt}

\newcommand{\figroot}{../Project/results/presentation/run_full_presentable/figures/analysis}
\newcommand{\statsroot}{../Project/results/presentation/run_full_presentable/figures/stats}
\newcommand{\reportfigroot}{../Project/results/presentation/run_full_presentable/figures/report}

\begin{document}

\pagenumbering{gobble}
\begin{titlepage}
\centering
{\LARGE Automated Negotiation Agents with NegMAS \par}
\vspace{0.4em}
{\Large Implementation, Strategy Comparison and Efficiency Analysis\par}
\vspace{1.5em}
{\large Adyant Srinivasan, Manan Gupta\par}
\vspace{0.6em}
{\normalsize University of California, Santa Barbara\par}
\vspace{0.6em}
{\normalsize ECE 270 - Game Theory\par}
\vspace{2em}
{\normalsize Date of Report: March 11, 2026\par}
\vfill
\begin{minipage}{0.95\linewidth}
\small
\end{minipage}
\end{titlepage}

\pagenumbering{roman}
\begin{abstract}
This project studies automated negotiation with NegMAS through a controlled buyer-seller procurement game and a multilateral extension. We implement and compare five negotiators: Boulware, Linear, Conceder, a custom SmartAspiration policy and a GeniusBOA baseline. The bilateral environment uses SAO bargaining on four issues (price, quantity, delivery time, service level), six preference families (A-F), three deadlines (12, 24, 36) and full ordered pairing across 40 seeds, yielding 18,000 negotiations. We evaluate agreement probability, rounds, individual utilities, social welfare, Pareto distance, Nash/Kalai/KS distances and efficiency losses (money left on table and welfare regret), with bootstrap confidence intervals and FDR-controlled pairwise permutation tests. We then test robustness through a multilateral procurement game (15,000 negotiations) and a random-scenario appendix tournament (600 negotiations). The results show strong deadline sensitivity, a consistent agreement-efficiency advantage for conceder behavior and a utility-agreement trade-off where the custom smart agent attains the highest average own utility but lower settlement rates. This report presents the model, implementation, empirical findings and practical lessons for strategy choice under different objectives.
\end{abstract}

\tableofcontents
\clearpage
\pagenumbering{arabic}

\section{Introduction}
\subsection{Background and Motivation}
Negotiation is a central interaction model in game theory because it combines private preferences, strategic behavior and constrained communication. In many engineering and computational settings, negotiation is also operationally important: supplier contracts, service-level agreements and procurement policies are often formed by iterative offers rather than one-shot auctions. As a result, it is useful to study not just theoretical bargaining solutions, but algorithmic negotiators running in explicit protocols with finite horizons.

Our project focuses on automated negotiating agents implemented in NegMAS. The key motivation is to bridge textbook bargaining concepts (Pareto efficiency, disagreement points, Nash/Kalai references, concession dynamics) with measurable computational behavior in repeated simulations. We intentionally use a game that is simple enough to explain clearly, but rich enough to reveal strategy trade-offs.

\subsection{Research Questions}
We organize the work around three questions:
\begin{itemize}[leftmargin=1.5em]
    \item \textbf{Q1:} How do standard time-based negotiating strategies compare in agreement rate, speed and efficiency under multiple deadlines?
    \item \textbf{Q2:} Does an opponent-aware aspiration strategy improve practical outcomes relative to standard built-in negotiators?
    \item \textbf{Q3:} Are bilateral findings stable when moving to a harder multilateral procurement setting?
\end{itemize}

\subsection{Contributions}
This report contributes:
\begin{enumerate}[leftmargin=1.5em]
    \item A complete, reproducible NegMAS implementation of bilateral and multilateral procurement-style negotiation experiments.
    \item A strategy benchmark over five negotiators, including a custom smart aspiration negotiator.
    \item A large-scale evaluation: 18,000 bilateral runs, 15,000 multilateral runs and a 600-negotiation appendix tournament.
    \item A statistically grounded analysis with bootstrap intervals, permutation testing, effect sizes and FDR correction.
\end{enumerate}

\section{Negotiation Game Model}
\subsection{Protocol}
The bilateral game is run with NegMAS \texttt{SAOMechanism} (Stacked Alternating Offers). Let \(o_t \in \Omega\) denote the offer at round \(t\), and let \(H_t = (o_1,\dots,o_t)\) be the offer history. Two agents alternate offers until one of three outcomes occurs:
\begin{enumerate}[leftmargin=1.5em]
    \item one side accepts an offer (agreement),
    \item the deadline is reached,
    \item the negotiation ends by time/penalty constraints.
\end{enumerate}
For each agent \(i\), the implemented acceptance rule is threshold-based:
\[
\text{accept at } t \iff u_i(o_t) \ge a_i(t) - \epsilon,
\]
where \(a_i(t)\) is the agent aspiration level and \(\epsilon\) is a small numerical tolerance.

In the full run, each bilateral session used \(n\_steps \in \{12,24,36\}\), \texttt{time\_limit=1.0s}, \texttt{pend=0.04} and \texttt{pend\_per\_second=0.02}. These settings intentionally make late rounds costly and reduce trivial always-agree behavior.

\subsection{Outcome Space}
The bilateral domain contains four issues:
\begin{itemize}[leftmargin=1.5em]
    \item \textbf{price}: 10 discrete levels,
    \item \textbf{quantity}: 10 discrete levels,
    \item \textbf{delivery\_time}: 10 discrete levels,
    \item \textbf{service\_level}: \{basic, standard, premium\}.
\end{itemize}
This structure allows conflicting monotonic preferences and categorical trade-offs.

\subsection{Utility Model and Profiles}
Each side has a private additive utility:
\[
u_i(o)=\sum_{k=1}^{K} w_{i,k} v_{i,k}(o_k), \quad K=4
\]
with nonlinear issue-value mappings \(v_{i,k}\) (linear/convex/concave depending on profile) and profile-specific weights \(w_{i,k}\). Utilities are scaled for comparability and include reserved values (\(0.35\) buyer and seller in the main study). If no agreement occurs, the implemented disagreement utility is the reserved-value point \((r_b,r_s)\).

Primary welfare and regret quantities are:
\[
W(o)=u_b(o)+u_s(o), \qquad
\text{regret}(o)=\frac{W^\star - W(o)}{W^\star},
\]
where \(W^\star\) is the maximum welfare over the computed Pareto frontier. Distance-based metrics are Euclidean in utility space, e.g.
\[
d_{\text{Pareto}}(o)=\min_{p \in \mathcal{P}} \lVert u(o)-p \rVert_2,
\]
for frontier \(\mathcal{P}\).

We evaluate six bilateral profile families:
\begin{itemize}[leftmargin=1.5em]
    \item \textbf{A}: balanced baseline,
    \item \textbf{B}: delivery-sensitive asymmetry,
    \item \textbf{C}: nonlinear price/quantity tension,
    \item \textbf{D}: service-dominant preferences,
    \item \textbf{E}: delivery-pressure asymmetry,
    \item \textbf{F}: strong specialization (quantity vs price emphasis).
\end{itemize}

\subsection{Evaluation Metrics}
For each run, we collect:
\begin{itemize}[leftmargin=1.5em]
    \item agreement indicator,
    \item rounds to agreement (and effective rounds in selected plots),
    \item buyer/seller utilities (\(u_b, u_s\)),
    \item welfare \(W=u_b+u_s\),
    \item distance to Pareto frontier,
    \item distances to Nash, Kalai and KS reference points,
    \item money left on table \(W^*_{\max}-W\),
    \item welfare regret \((W^*_{\max}-W)/W^*_{\max}\).
\end{itemize}

\section{Negotiation Strategies}
\subsection{Baseline Strategies}
We benchmark three standard time-based negotiators:
\begin{itemize}[leftmargin=1.5em]
    \item \texttt{BoulwareTBNegotiator}: late concession, typically hard bargaining.
    \item \texttt{LinearTBNegotiator}: approximately linear concession.
    \item \texttt{ConcederTBNegotiator}: early concession to increase settlement probability.
\end{itemize}

\subsection{Additional Baseline}
\texttt{GeniusBOABaselineNegotiator} is implemented as a BOA-inspired aspiration baseline inside NegMAS (aspiration-style concession and acceptance behavior).

\subsection{Custom Strategy: SmartAspirationNegotiator}
Our custom strategy keeps an aspiration threshold over time but modifies offer generation using opponent behavior. The implementation in \texttt{src/negotiators.py} is a two-phase aspiration policy with partner-profile proximity filtering.

Let \(t \in [0,1]\) be relative time in the negotiation. The normalized aspiration utility is piecewise:
\[
\alpha(t)=
\begin{cases}
\alpha_{\text{early}}(t), & t \le \tau, \\
\alpha_{\text{early}}(\tau)\,\alpha_{\text{late}}\!\left(\frac{t-\tau}{1-\tau}\right), & t>\tau,
\end{cases}
\]
where \(\tau=0.8\) in our run configuration, \(\alpha_{\text{early}}\) is Boulware-style, and \(\alpha_{\text{late}}\) is Linear-style. The utility threshold is then:
\[
a(t)=u_{\min} + (u_{\max}-u_{\min})\alpha(t).
\]

Define the feasible offer set:
\[
\mathcal{C}_t = \{o \in \Omega : u_i(o) \ge a(t)\}.
\]
The strategy estimates a lightweight opponent profile from observed offers \(H_t\):
\begin{itemize}[leftmargin=1.5em]
    \item for numeric issue \(k\): target \(\mu_k=\text{mean}(o_{1,k},\dots,o_{t,k})\), weight \(\omega_k \propto 1/(\text{span}_k+10^{-3})\);
    \item for categorical issue \(k\): target is the modal value, with \(\omega_k \propto 1/|\text{unique values}|\).
\end{itemize}
Weights are normalized to \(\sum_k \omega_k=1\). For candidate \(o\), distance to partner profile is:
\[
D(o)=\sum_k \omega_k\,\delta_k(o_k,\mu_k),
\]
with normalized absolute distance for numeric issues and \(0/1\) mismatch for categorical issues. Proposal is:
\[
o_t^{\text{propose}}=\arg\min_{o \in \mathcal{C}_t} D(o).
\]
Acceptance still follows the same aspiration threshold \(u_i(o_t)\ge a(t)-\epsilon\).

\section{Implementation Details}
\subsection{Code Architecture}
The project implementation is organized into small modules so that experiment logic, game definition and analytics are separated:
\begin{itemize}[leftmargin=1.5em]
    \item \texttt{src/domain.py}: outcome space construction, profile-specific utility generators and SAO mechanism setup.
    \item \texttt{src/negotiators.py}: baseline wrappers and custom \texttt{SmartAspirationNegotiator}.
    \item \texttt{src/metrics.py}: welfare, frontier distances, bootstrap confidence intervals, permutation tests and FDR correction.
    \item \texttt{src/run\_experiments.py}: bilateral matrix runner and CSV artifact export.
    \item \texttt{src/run\_multilateral\_experiments.py}: one-buyer/two-supplier extension.
    \item \texttt{src/appendix\_tournament.py}: random-scenario robustness tournament.
    \item \texttt{scripts/plot\_stats.py} and \texttt{scripts/plot\_analysis.py}: figure generation from saved CSV results.
\end{itemize}

\subsection{Execution Pipeline and Validation}
The final bundle \texttt{run\_full\_presentable} was produced by one sequential pipeline run with status \texttt{success}. Before large simulation, the workflow executes a smoke test and \texttt{pytest} for basic correctness checks, then runs bilateral, appendix and multilateral stages, and finally generates all plots.

\begin{table}[H]
\centering
\caption{Sequential execution stages from \texttt{manifest.json}}
\small
\begin{tabular}{lrr}
\toprule
Stage & Duration (s) & Share of total \\
\midrule
Smoke test & 0.31 & 0.02\% \\
Pytest & 1.32 & 0.08\% \\
Main experiments & 1129.31 & 72.08\% \\
Plot stats & 38.04 & 2.43\% \\
Appendix tournament & 28.11 & 1.79\% \\
Multilateral experiments & 361.44 & 23.07\% \\
Plot analysis & 8.16 & 0.52\% \\
\midrule
Total & 1566.69 & 100.00\% \\
\bottomrule
\end{tabular}
\end{table}

\subsection{Data Contracts}
The bilateral raw file \texttt{main/main\_runs.csv} stores one row per negotiation with run identity, profile, deadline, ordered strategy pairing, agreement flag/outcome, round count, utilities and all efficiency metrics. Aggregated files \texttt{main/main\_summary.csv}, \texttt{main/main\_strategy\_summary.csv} and \texttt{main/main\_pairwise\_tests.csv} support report tables and significance claims. Equivalent raw/summary outputs exist for multilateral and appendix stages.

\section{Experimental Design and Reproducibility}
\subsection{Bilateral Main Matrix}
\begin{table}[H]
\centering
\caption{Bilateral design used in the full run}
\begin{tabular}{ll}
\toprule
\textbf{Component} & \textbf{Setting} \\
\midrule
Profiles & A, B, C, D, E, F \\
Deadlines & 12, 24, 36 \\
Strategies & 5 (all ordered pairs, self-play included) \\
Seeds per condition & 40 \\
Runs & \(6 \times 3 \times 5 \times 5 \times 40 = 18{,}000\) \\
Penalty settings & \texttt{pend=0.04}, \texttt{pend\_per\_second=0.02} \\
Time cap per negotiation & \texttt{time\_limit=1.0s} \\
\bottomrule
\end{tabular}
\end{table}

\subsection{Multilateral Extension}
The multilateral setup models one buyer and two suppliers (3 agents total) over profile families \(M1, M2, M3\), \(n\_steps \in \{20,40\}\) and all ordered strategy vectors:
\[
3 \times 2 \times 5^3 \times 20 = 15{,}000 \text{ runs}
\]

\subsection{Appendix Tournament}
We additionally run a randomized robustness appendix using NegMAS cartesian tournament:
\begin{itemize}[leftmargin=1.5em]
    \item 6 random scenarios,
    \item 2 repetitions,
    \item 5 competitors,
    \item 600 negotiations total.
\end{itemize}

\subsection{Statistical Analysis}
For aggregated results we report mean, standard deviation and 95\% confidence intervals. The analytic CI reported in summaries is:
\[
\bar{x} \pm 1.96\frac{s}{\sqrt{n}},
\]
and bootstrap CIs are also computed with 2000 resamples for robustness.

Pairwise strategy comparisons are performed using two-sided permutation tests (\(n=2000\) permutations) on mean differences. Multiple testing is controlled per context/metric using Benjamini-Hochberg FDR:
\[
p^{\text{FDR}}_{(i)} = \min_{j \ge i}\left(\frac{m}{j}p_{(j)}\right),
\]
where \(p_{(j)}\) are sorted raw p-values and \(m\) is the number of tests in that family. We also report Cohen's \(d\) as standardized effect size.

\section{Results}
\subsection{Overall Outcomes}
\begin{table}[H]
\centering
\caption{High-level outcomes from \texttt{run\_full\_presentable}}
\begin{tabular}{lcc}
\toprule
\textbf{Experiment} & \textbf{Runs} & \textbf{Agreement Rate} \\
\midrule
Bilateral main & 18,000 & 0.4201 \\
Multilateral extension & 15,000 & 0.3979 \\
Appendix tournament & 600 & score-ranked \\
\bottomrule
\end{tabular}
\end{table}

From \texttt{manifest.json}, this run completed end-to-end in 1566.69 seconds (\(\approx\)26.1 minutes), with bilateral simulation dominating runtime (72.08\%) and multilateral simulation second (23.07\%). The pipeline status is \texttt{success}, and both smoke test and pytest stages pass before experiment execution.

The bilateral agreement rate of 0.4201 and multilateral agreement rate of 0.3979 indicate that both environments are non-trivial under this penalty/deadline configuration. This is desirable for comparison studies: weak strategies fail often enough to be distinguishable.

\subsection{Deadline and Profile Effects}
Agreement drops sharply with larger \(n\_steps\) in both bilateral and multilateral setups:
\begin{itemize}[leftmargin=1.5em]
    \item Bilateral: 12-step 0.6512, 24-step 0.3815, 36-step 0.2277.
    \item Multilateral: 20-step 0.5156, 40-step 0.2801.
\end{itemize}
Given positive reservation values and penalties over negotiation progression, this pattern is consistent with stricter effective acceptance conditions as negotiations drag on.

\begin{figure}[H]
\centering
\includegraphics[width=\linewidth]{\reportfigroot/figure1_agreement_clean.png}
\caption{Revised agreement view for readability. Panel A shows agreement rate versus deadline for each profile (A--F), aggregated over strategy pairings. Panel B shows ordered buyer-seller pairing agreement rates, aggregated over profiles and deadlines, with short strategy labels.}
\label{fig:agree-pair}
\end{figure}

\begin{table}[H]
\centering
\caption{Bilateral profile-deadline summary (agreement, welfare, Pareto distance, regret)}
\small
\begin{tabular}{cccccc}
\toprule
Profile & Steps & Agreement & Welfare & Pareto Dist. & Regret \\
\midrule
A & 12 & 0.6540 & 0.6559 & 0.2433 & 0.3441 \\
A & 24 & 0.3620 & 0.3654 & 0.4487 & 0.6346 \\
A & 36 & 0.2120 & 0.2162 & 0.5542 & 0.7838 \\
B & 12 & 0.6590 & 0.7292 & 0.3170 & 0.4132 \\
B & 24 & 0.4160 & 0.4793 & 0.4978 & 0.6143 \\
B & 36 & 0.2550 & 0.2933 & 0.6316 & 0.7640 \\
C & 12 & 0.6530 & 0.7132 & 0.3631 & 0.4442 \\
C & 24 & 0.3870 & 0.4397 & 0.5646 & 0.6574 \\
C & 36 & 0.2370 & 0.2722 & 0.6900 & 0.7879 \\
D & 12 & 0.6360 & 0.5879 & 0.2966 & 0.4121 \\
D & 24 & 0.3480 & 0.3256 & 0.4854 & 0.6744 \\
D & 36 & 0.2110 & 0.2012 & 0.5743 & 0.7988 \\
E & 12 & 0.6430 & 0.6732 & 0.3067 & 0.4520 \\
E & 24 & 0.3720 & 0.3970 & 0.4841 & 0.6768 \\
E & 36 & 0.2020 & 0.2082 & 0.6072 & 0.8305 \\
F & 12 & 0.6620 & 0.8121 & 0.3799 & 0.4630 \\
F & 24 & 0.4040 & 0.5008 & 0.5817 & 0.6689 \\
F & 36 & 0.2490 & 0.3088 & 0.7141 & 0.7958 \\
\bottomrule
\end{tabular}
\end{table}

\noindent\textbf{Mathematical interpretation.} As deadlines increase from 12 to 36 steps, mean agreement decreases while both \(d_{\text{Pareto}}\) and regret increase. Through
\[
\text{regret}=\frac{W^\star-W}{W^\star},
\]
this implies negotiated welfare \(W\) falls farther below frontier-optimal welfare \(W^\star\) in longer-horizon runs under the chosen penalty/reservation configuration. In other words, the additional rounds do not improve exploration enough to offset the stricter effective acceptability induced by time progression.

\subsection{Bilateral Strategy Comparison}
Table-level outcomes show a clear pattern:
\begin{itemize}[leftmargin=1.5em]
    \item \texttt{Conceder} maximizes agreement and minimizes regret.
    \item \texttt{SmartAspiration} maximizes mean own utility but with lower agreement.
    \item \texttt{Linear} sits between fast agreement and hard bargaining.
    \item \texttt{Boulware}/\texttt{GeniusBOA} settle less frequently in this configuration.
\end{itemize}

\begin{table}[H]
\centering
\caption{Role-aggregated bilateral strategy means with 95\% CI (\(n=7200\) role-observations per strategy)}
\small
\begin{tabular}{lccccccc}
\toprule
Strategy & Own Util. & Agree. & Agree CI & Welfare & Welfare CI & Pareto Dist. & Regret \\
\midrule
SmartAspiration & 0.2753 & 0.3981 & \(\pm 0.0113\) & 0.4523 & \(\pm 0.0130\) & 0.4893 & 0.6261 \\
Boulware & 0.2324 & 0.3414 & \(\pm 0.0110\) & 0.3717 & \(\pm 0.0120\) & 0.5464 & 0.6925 \\
Linear & 0.2315 & 0.4603 & \(\pm 0.0115\) & 0.5003 & \(\pm 0.0127\) & 0.4547 & 0.5861 \\
GeniusBOA & 0.2278 & 0.3401 & \(\pm 0.0109\) & 0.3644 & \(\pm 0.0118\) & 0.5507 & 0.6978 \\
Conceder & 0.1690 & 0.5607 & \(\pm 0.0115\) & 0.5833 & \(\pm 0.0121\) & 0.3867 & 0.5130 \\
\bottomrule
\end{tabular}
\end{table}

The uncertainty ranges are small because each strategy aggregates 7,200 role-observations. Even with uncertainty included, the ordering remains stable for the core trade-off: \texttt{Conceder} has the highest agreement and lowest regret, while \texttt{SmartAspiration} has the highest mean own utility. This means the central conclusion is not an artifact of noisy small-sample means.

\begin{table}[H]
\centering
\caption{Top ordered bilateral pairings by welfare}
\small
\begin{tabular}{llcccc}
\toprule
Buyer & Seller & Agree. & Welfare & Pareto Dist. & Regret \\
\midrule
Conceder & Conceder & 0.8153 & 0.7961 & 0.2736 & 0.3357 \\
Linear & Conceder & 0.6819 & 0.7568 & 0.2846 & 0.3697 \\
Conceder & Linear & 0.6681 & 0.7054 & 0.2880 & 0.4115 \\
SmartAspiration & Conceder & 0.5514 & 0.5790 & 0.3736 & 0.5146 \\
Linear & Linear & 0.5000 & 0.5446 & 0.4386 & 0.5589 \\
SmartAspiration & Linear & 0.4764 & 0.5390 & 0.4311 & 0.5551 \\
\bottomrule
\end{tabular}
\end{table}

\begin{figure}[H]
\centering
\begin{subfigure}{0.49\linewidth}
\centering
\includegraphics[width=\linewidth]{\figroot/analysis_heatmap_agreement.png}
\caption{Agreement}
\end{subfigure}
\hfill
\begin{subfigure}{0.49\linewidth}
\centering
\includegraphics[width=\linewidth]{\figroot/analysis_heatmap_welfare.png}
\caption{Welfare}
\end{subfigure}
\caption{Ordered matchup heatmaps (buyer strategy on rows, seller strategy on columns).}
\label{fig:heatmap-main}
\end{figure}

\begin{figure}[H]
\centering
\begin{subfigure}{0.49\linewidth}
\centering
\includegraphics[width=\linewidth]{\figroot/analysis_heatmap_pareto_distance.png}
\caption{Pareto distance}
\end{subfigure}
\hfill
\begin{subfigure}{0.49\linewidth}
\centering
\includegraphics[width=\linewidth]{\figroot/analysis_heatmap_effective_rounds.png}
\caption{Effective rounds (failures at deadline)}
\end{subfigure}
\caption{Efficiency and duration heatmaps by ordered strategy pairing.}
\label{fig:heatmap-extra}
\end{figure}

\subsection{Efficiency and Reference-Point Analysis}
Beyond agreement and utility, distance-to-reference metrics show meaningful variation across profiles and strategies. On average across all bilateral runs:
\begin{itemize}[leftmargin=1.5em]
    \item mean Pareto distance \(=0.4856\),
    \item mean Nash distance \(=0.6216\),
    \item mean KS distance \(=0.6156\),
    \item mean welfare regret \(=0.6231\).
\end{itemize}
These values confirm that a large fraction of outcomes remain efficiency-suboptimal under realistic deadline penalties.

\begin{figure}[H]
\centering
\begin{subfigure}{0.49\linewidth}
\centering
\includegraphics[width=\linewidth]{\figroot/analysis_reference_distances.png}
\caption{Nash, Kalai, KS distances}
\end{subfigure}
\hfill
\begin{subfigure}{0.49\linewidth}
\centering
\includegraphics[width=\linewidth]{\figroot/analysis_efficiency_metrics.png}
\caption{Money left on table, regret}
\end{subfigure}
\caption{Reference-point and efficiency diagnostics from the bilateral full run.}
\label{fig:ref-eff}
\end{figure}

\subsection{Significance Testing}
The pairwise analysis file contains 3,960 comparisons across contexts and metrics. Using the criterion \(p^{\text{FDR}}<0.05\), 2,008 tests are significant (50.71\%). The largest counts of significant pairwise differences appear in:
\begin{itemize}[leftmargin=1.5em]
    \item rounds (247),
    \item money left on table (186),
    \item Pareto distance (185),
    \item welfare regret (181),
    \item agreement (179).
\end{itemize}
This means major conclusions are not only visible in means but also statistically stable across many profile/deadline contexts.

\subsection{Multilateral Results}
In the one-buyer/two-supplier game:
\begin{itemize}[leftmargin=1.5em]
    \item \texttt{Conceder} as buyer has the best agreement (0.5297), highest welfare (1.0638) and lowest welfare regret (0.5206).
    \item \texttt{SmartAspiration} has lower agreement (0.3227) and higher regret (0.7099).
    \item Deadline effect persists strongly: 20-step agreement 0.5156 vs 40-step agreement 0.2801.
\end{itemize}

\begin{table}[H]
\centering
\caption{Multilateral buyer-strategy means with 95\% CI (\(n=3000\) per buyer strategy)}
\small
\begin{tabular}{lcccccc}
\toprule
Buyer strategy & Agreement & Agree CI & Buyer Util. & Welfare & Welfare CI & Regret \\
\midrule
Conceder & 0.5297 & \(\pm 0.0179\) & 0.2278 & 1.0638 & \(\pm 0.0363\) & 0.5206 \\
Linear & 0.4357 & \(\pm 0.0177\) & 0.2332 & 0.8753 & \(\pm 0.0360\) & 0.6088 \\
Boulware & 0.3507 & \(\pm 0.0171\) & 0.2426 & 0.6792 & \(\pm 0.0333\) & 0.6958 \\
GeniusBOA & 0.3507 & \(\pm 0.0171\) & 0.2399 & 0.6647 & \(\pm 0.0327\) & 0.7019 \\
SmartAspiration & 0.3227 & \(\pm 0.0167\) & 0.2262 & 0.6509 & \(\pm 0.0340\) & 0.7099 \\
\bottomrule
\end{tabular}
\end{table}

\begin{figure}[H]
\centering
\begin{subfigure}{0.49\linewidth}
\centering
\includegraphics[width=\linewidth]{\figroot/analysis_multilateral_agreement_by_buyer_strategy.png}
\caption{Agreement by buyer strategy}
\end{subfigure}
\hfill
\begin{subfigure}{0.49\linewidth}
\centering
\includegraphics[width=\linewidth]{\figroot/analysis_multilateral_welfare_regret_by_buyer_strategy.png}
\caption{Welfare regret by buyer strategy}
\end{subfigure}
\caption{Multilateral procurement outcomes from the full run.}
\label{fig:multi}
\end{figure}

\subsection{Appendix Tournament}
The appendix score ranking by NegMAS-provided \texttt{advantage} is:
\begin{enumerate}[leftmargin=1.5em]
    \item SmartAspiration (0.7070)
    \item Boulware (0.6975)
    \item GeniusBOA (0.6924)
    \item Linear (0.4789)
    \item Conceder (0.2990)
\end{enumerate}
This differs from fixed-domain welfare/agreement rankings and highlights that ``best strategy'' depends on objective, scenario generation and tournament scoring definitions. In this report, appendix ranking is treated strictly as a robustness supplement, not as a replacement for the fixed-domain primary claims.

\section{Discussion}
The results support several robust interpretations.

\textbf{First}, the project clearly exhibits a utility-agreement trade-off. Conceder behavior increases settlement probability and reduces regret but yields lower own utility. SmartAspiration behaves in the opposite direction: better individual utility but weaker agreement robustness.

\textbf{Second}, deadline and penalty design strongly shape empirical outcomes. The decline in agreement with larger \(n\_steps\) may appear counter-intuitive in isolation, but it is coherent here because reservation values and progression penalties effectively tighten acceptance conditions over longer traces.

\textbf{Third}, ordered role analysis matters. Buyer-only or seller-only aggregation can hide interaction effects. Pairing heatmaps reveal that outcomes are jointly determined by both roles and asymmetries are not negligible.

\textbf{Fourth}, bilateral conclusions partially transfer to multilateral settings. The same broad pattern (Conceder stronger on agreement/efficiency) remains visible, but magnitudes and ranking gaps change due to added coordination complexity.

\section{Limitations}
\begin{itemize}[leftmargin=1.5em]
    \item \textbf{Model scope:} utility profiles are synthetic families, not fit from real procurement data.
    \item \textbf{Protocol scope:} all main results use SAO; other mechanisms may reorder rankings.
    \item \textbf{Time cap effects:} \texttt{time\_limit=1.0s} improves reliability but may truncate some traces.
    \item \textbf{Baseline interpretation:} \texttt{GeniusBOABaselineNegotiator} is BOA-inspired within NegMAS rather than an external GENIUS bridge replication.
    \item \textbf{Metric dependence:} appendix tournament ``advantage'' is a NegMAS score and is not identical to welfare-optimality.
\end{itemize}

\section{Conclusion}
We implemented and evaluated a complete automated negotiation study aligned with the course proposal but substantially expanded in scale and rigor. The final evidence from \texttt{run\_full\_presentable} shows that:
\begin{enumerate}[leftmargin=1.5em]
    \item concession-heavy strategies improve agreement and efficiency,
    \item opponent-aware aspiration can improve own utility at agreement-risk cost,
    \item strategy quality is objective-dependent and role-dependent,
    \item conclusions remain directionally stable in a multilateral extension.
\end{enumerate}

For practical deployment, the key recommendation is to select strategy by objective: if reliability and social efficiency are primary, use conceder-like policies; if individual payoff is primary and failed talks are acceptable, use smarter aspiration-based bargaining with stronger opponent modeling.

\section*{References}
\begin{enumerate}[leftmargin=1.5em]
    \item Y. M. A. Ezzat, ``NegMAS documentation,'' v0.15.1.post1, 2026. \url{https://negmas.readthedocs.io}
    \item J. F. Nash, ``The Bargaining Problem,'' \emph{Econometrica}, vol. 18, no. 2, pp. 155--162, 1950.
    \item E. Kalai and M. Smorodinsky, ``Other Solutions to Nash's Bargaining Problem,'' \emph{Econometrica}, vol. 43, no. 3, pp. 513--518, 1975.
    \item P. Faratin, C. Sierra and N. R. Jennings, ``Negotiation decision functions for autonomous agents,'' \emph{Robotics and Autonomous Systems}, vol. 24, no. 3--4, pp. 159--182, 1998.
    \item T. Baarslag, M. Kaisers, E. H. Gerding, C. M. Jonker and J. Gratch, ``Evaluating practical negotiating agents: Results and analysis of the 2011--2016 automated negotiating agents competition,'' \emph{Artificial Intelligence}, vol. 252, pp. 66--92, 2017.
\end{enumerate}

\end{document}
